\section{Zeta-regularized determinant}

We recall the notion of the zeta-regularized determinant.
Let $\Theta:V\rightarrow{V}$ be a linear operator acting on a complex vector space $V$ of countable dimension.
We assume that $V$ is the direct sum of finite-dimensional $\Theta$-invariant sub-spaces.
Let $\mathrm{Sp}(\Theta)$ be the set of eigenvalues of $\Theta$. 
The spectral zeta function associated with the operator $\Theta$ is defined by the analytic continuation of Dirichlet series
\begin{equation*}
  \zeta_{\Theta}(s)=\sum_{\lambda\ne0\in\mathrm{Sp}(\Theta)}\lambda^{-s}\text{ with }\lambda^{-s}=|\lambda|^{-s}e^{-is(\mathrm{Arg}\lambda)}, -\pi<\mathrm{Arg}\lambda\le\pi.
\end{equation*}
We assume that the Dirichlet series converges absolutely on some right-half plane and has an analytic continuation to the half plane $\mathrm{Re}(s)>-\epsilon$ for some $\epsilon>0$ which is holomorphic at $s=0$.
Under these conditions, we define a zeta-regularized determinant by
\begin{equation*}
    \mathrm{det}_{\infty}(\Theta|V):=\exp\left(-\partial_{s}\zeta_{\Theta}(0)\right).
\end{equation*}
